% 原文第35页；公式保留原稿矢量内容。
\zhbox{54}{64}{505}{16}{13}{\textbf{应用信息}}
\zhbox{54}{83}{235}{79}{10}{\textbf{输出电压设置}\\默认输出电压 $V_{OUT}$ 由 FB 电阻分压器设置。默认参考电压 $V_{REF}$ 为 0.5V。分压器下端电阻通常为 1k$\Omega$ 至 50k$\Omega$；上端电阻 R1 可由式（6）计算：}
\zhbox{54}{207}{235}{14}{9.5}{也可通过 I\textsuperscript{2}C 接口选择 $V_{REF}$，以设置 $V_{OUT}$。}
\zhbox{54}{226}{235}{67}{10}{\textbf{电感选择}\\电感的选择取决于工作模式（降压、升压或升降压模式）。降压模式的电感量（$I_{L\_BUCK}$）可由式（7）计算：}
\zhbox{54}{340}{235}{57}{9.5}{其中，$\Delta I_L$ 为电感纹波电流峰峰值。\\选择 $\Delta I_L$ 约为最大负载电流 $I_{LOAD\_MAX}$ 的 30\% 至 50\%。}
\zhbox{54}{403}{235}{63}{9.5}{升压模式下，应根据 $\Delta I_L$ 确定电感。选择 $\Delta I_L$ 约为最大输入电流 $I_{IN\_MAX}$ 的 30\% 至 50\%。升压模式的电感量（$I_{L\_BOOST}$）可由式（8）计算：}
\zhbox{54}{515}{235}{25}{9.5}{最大输入电流 $I_{IN\_MAX}$ 可由式（9）计算：}
\zhbox{54}{589}{235}{25}{9.5}{其中，$\Delta I_L$ 约为 $I_{IN\_MAX}$ 的 30\% 至 50\%，$\eta$ 为效率。}
\zhbox{54}{620}{235}{78}{9.5}{较大的电感可减小 $\Delta I_L$，但会占用更多空间，并将右半平面零点移向较低频率，从而降低转换器可达到的带宽。应根据具体应用要求选择电感。}
\zhbox{324}{83}{235}{193}{10}{\textbf{输入电容（$C_{IN}$）选择}\\降压模式下，转换器的输入电流 $I_{IN}$ 不连续，需要电容在维持直流输入电压的同时，向转换器提供交流电流。（升压模式下，转换器的输入电流连续。）为获得最佳性能，建议使用陶瓷电容，并尽量靠近 VIN 放置。建议使用 X5R 或 X7R 介质的陶瓷电容，其 ESR 低、温度系数小。\\选择纹波电流额定值大于最大输入纹波电流 $\Delta I_{IN}$ 的电容。降压模式的输入纹波电流 $I_{CIN\_RMS}$ 可由式（10）估算：}
\zhbox{324}{326}{235}{38}{9.5}{降压模式下，最不利工况发生在 $V_{IN}=2\times V_{OUT}$ 时，可由式（11）计算：}
\zhbox{324}{410}{235}{39}{9.5}{为简化选择，输入电容 $C_{IN}$ 的 RMS 电流额定值应大于 $I_{LOAD\_MAX}$ 的一半。}
\zhbox{324}{454}{235}{39}{9.5}{$C_{IN}$ 决定转换器的输入电压纹波 $\Delta V_{IN}$。若系统对 $\Delta V_{IN}$ 有要求，应选择符合系统规格的 $C_{IN}$。}
\zhbox{324}{498}{235}{26}{9.5}{降压模式下，$\Delta V_{IN}$ 可由式（12）估算：}
\zhbox{324}{572}{235}{26}{9.5}{最不利工况发生在 $V_{IN}=2\times V_{OUT}$ 时，可由式（13）计算：}
\zhbox{324}{646}{235}{43}{9.5}{\textbf{输出电容（$C_{OUT}$）选择}\\升压模式下，转换器的输出电压（$I_{OUT}$）不连续。输出}
